%%% FIELD uniform-optimum-proof
By \texttt{ass:supply-range} and \texttt{def:uniform-dosage},
\(u=\min\{L/p,1/2\}\in(0,1/2]\), and \(d^\star=u\mathbf 1_p\) is an interior point of the cube and belongs to \(\mathcal D_L\).  Indeed, if \(L<p/2\), then \(\sum_i d_i^\star=pu=L\), whereas if \(L\ge p/2\), then \(\sum_i d_i^\star=p/2\le L\).

We first establish the curvature certificate on the open cube.  For \(0<v<1\), put \(m=2v-1\).  The definition \texttt{def:coordinate-factor} gives
\[
 r(v)=\frac{1+m^2}{1-m^2}\ge 1,
\qquad
 r'(v)=-\frac{1-2v}{2v^2(1-v)^2},
\]
where \(|m|<1\) makes the first denominator positive.  Direct differentiation also gives the exact log-curvature identity
\[
 (\log r)''(v)
 =16\frac{1+3m^4}{(1-m^4)^2}
 =16+16\frac{m^4(5-m^4)}{(1-m^4)^2}
 \ge 16. \tag{1}
\]
The last inequality follows from \(0\le m^4<1\), so every factor and denominator in the remainder is nonnegative or positive as appropriate.

By the already established \texttt{lem:trace-identity}, on \((0,1)^p\) one has
\[
 \Phi_k(d)=\sum_{T\in\mathcal I_k} f_T(d),
 \qquad
 f_T(d):=\prod_{i\in T}r(d_i).
\]
For a nonempty \(T\), define the vector \(a_T(d)\in\mathbb R^p\) and diagonal matrix \(D_T(d)\in\mathbb R^{p\times p}\) by
\[
 (a_T(d))_i=\mathbf 1\{i\in T\}(\log r)'(d_i),
 \qquad
 (D_T(d))_{ii}=\mathbf 1\{i\in T\}(\log r)''(d_i).
\]
Twice differentiating the finite product yields
\[
 \nabla^2 f_T(d)
 =f_T(d)\bigl(D_T(d)+a_T(d)a_T(d)^{\mathsf T}\bigr).
\]
The matrix \(a_Ta_T^{\mathsf T}\) is positive semidefinite.  Moreover, \(f_T(d)\ge1\) by the first display and every active diagonal entry of \(D_T\) is at least \(16\) by (1).  Hence
\[
 \nabla^2 f_T(d)\succeq
 16\operatorname{diag}\bigl((\mathbf 1\{i\in T\})_{i=1}^p\bigr). \tag{2}
\]
For a fixed coordinate \(i\), the number of nonempty sets \(T\in\mathcal I_k\) containing \(i\) is, by \texttt{def:index-family} and \texttt{def:curvature-multiplicity},
\[
 \sum_{h=0}^{k-1}\binom{p-1}{h}=C_{p,k};
\]
one chooses the other \(h=|T|-1\) elements from the remaining \(p-1\) coordinates.  Summing (2), while noting that the empty-set summand is constant, proves
\[
 \nabla^2\Phi_k(d)\succeq16C_{p,k}I_p
 \qquad(d\in(0,1)^p). \tag{3}
\]

Let \(d\in\mathcal D_L\cap(0,1)^p\), set \(\delta=d-d^\star\), and apply the one-dimensional fundamental theorem of calculus twice to \(g(t)=\Phi_k(d^\star+t\delta)\).  The whole segment lies in the open cube, and (3) gives
\[
\begin{aligned}
 \Phi_k(d)-\Phi_k(d^\star)
 &=\nabla\Phi_k(d^\star)^{\mathsf T}\delta
   +\int_0^1(1-t)\,\delta^{\mathsf T}
       \nabla^2\Phi_k(d^\star+t\delta)\delta\,dt\\
 &\ge \nabla\Phi_k(d^\star)^{\mathsf T}\delta
   +8C_{p,k}\|\delta\|_2^2. \tag{4}
\end{aligned}
\]
It remains to verify the sign of the first-order term rather than discard it without justification.  Differentiating the trace identity and using the fact that all coordinates of \(d^\star\) equal \(u\) gives, for every \(i\),
\[
 \frac{\partial\Phi_k}{\partial d_i}(d^\star)
 =r'(u)\sum_{h=0}^{k-1}\binom{p-1}{h}r(u)^h
 =:\gamma_u. \tag{5}
\]
If \(u=1/2\), the displayed formula for \(r'\) gives \(\gamma_u=0\).  If \(u=L/p<1/2\), then \(r'(u)<0\), while the sum in (5) is strictly positive because its \(h=0\) term is one; thus \(\gamma_u<0\).  Feasibility from \texttt{ass:expected-supply} gives
\[
 \sum_{i=1}^p\delta_i=\sum_{i=1}^p d_i-pu\le L-L=0,
\]
and consequently
\(\nabla\Phi_k(d^\star)^{\mathsf T}\delta
=\gamma_u\sum_i\delta_i\ge0\).  In either supply regime, (4) therefore proves
\[
 \Phi_k(d)-\Phi_k(d^\star)
 \ge8C_{p,k}\|d-d^\star\|_2^2. \tag{6}
\]
If \(d\) is on the boundary of \([0,1]^p\), at least one coordinate is zero or one, and \texttt{lem:trace-identity} gives \(\Phi_k(d)=+\infty\); hence (6) remains true in the extended-real sense.  Since \(C_{p,k}\ge1\), (6) is strict for every \(d\ne d^\star\), proving that \(d^\star\) is the unique minimizer over \(\mathcal D_L\).

Finally, substitution of the uniform dosage into the trace identity, followed by grouping sets according to their cardinality, gives the claimed exact value:
\[
 \Phi^\star=\Phi_k(d^\star)
 =\sum_{T\in\mathcal I_k}r(u)^{|T|}
 =\sum_{j=0}^k\binom pj r(u)^j. \tag{7}
\]

It remains to prove sharpness when \(L\ge p/2\).  Then \(d^\star=(1/2)\mathbf 1_p\).  At this point every \(f_T=1\), every \(a_T=0\), and every active diagonal entry of \(D_T\) equals \(16\) by (1).  The same counting used in (3) therefore gives the equality
\[
 \nabla^2\Phi_k(d^\star)=16C_{p,k}I_p. \tag{8}
\]
The assumption \(p\ge2\), inherited through \texttt{lem:trace-identity}, permits the nonzero direction \(h=e_1-e_2\).  For all sufficiently small \(|t|\), the point \(d(t)=d^\star+th\) lies in \((0,1)^p\), has coordinate sum \(p/2\le L\), and hence is feasible.  Applying the same twice-integrated identity as in (4), using \(\nabla\Phi_k(d^\star)=0\), continuity of the finite-sum Hessian, and (8), yields
\[
 \lim_{t\to0}
 \frac{\Phi_k(d(t))-\Phi_k(d^\star)}
      {\|d(t)-d^\star\|_2^2}
 =\frac{\frac12 h^{\mathsf T}(16C_{p,k}I_p)h}{\|h\|_2^2}
 =8C_{p,k}. \tag{9}
\]
Thus any larger coefficient in a global inequality of the form (6) would fail along this feasible path.  Together with (6), this proves that \(8C_{p,k}\) is the best global quadratic certificate when \(L\ge p/2\).
%%% FIELD additive-reduction-proof
Suppose \(k=1\).  By \texttt{def:index-family}, the retained family is
\(\mathcal I_1=\{\varnothing,\{1\},\ldots,\{p\}\}\).  Therefore the already established \texttt{lem:trace-identity} gives, for every interior feasible dosage,
\[
 \Phi_1(d)
 =\prod_{i\in\varnothing}r(d_i)
  +\sum_{i=1}^p\prod_{\ell\in\{i\}}r(d_\ell)
 =1+\sum_{i=1}^p r(d_i). \tag{1}
\]
At a boundary dosage the same trace lemma assigns \(\Phi_1(d)=+\infty\), consistently with the singular-design convention in \texttt{def:population-criterion}.

Applying \texttt{thm:uniform-optimum} with \(k=1\) shows that the unique optimizer is the uniform limited-supply dosage
\[
 d_i=u=\min\{L/p,1/2\}\qquad(1\le i\le p).
\]
This is precisely the additive product-Bernoulli allocation asserted in the proposition.  If \(L\ge p/2\), then \(u=1/2\), and \texttt{def:coordinate-factor} gives
\[
 r(1/2)=\frac{1}{2(1/2)(1/2)}-1=1.
\]
Substituting into (1), or equivalently into the exact value in \texttt{thm:uniform-optimum}, yields \(\Phi^\star=1+p\), as required.
